% ════════════════════════════════════════════════════════════════════
%  front_pages/abstract.tex
%  Abstract Page
% ════════════════════════════════════════════════════════════════════

\chapter*{Abstract}
\addcontentsline{toc}{chapter}{Abstract}
\thispagestyle{plain}
\vspace{0.3cm}

\begin{onehalfspace}
\noindent
Non-small cell lung cancer (NSCLC) accounts for approximately
85\% of all lung cancer cases globally, with epidermal growth
factor receptor (EGFR) mutations representing one of the most
clinically significant oncogenic drivers. EGFR-targeted tyrosine
kinase inhibitors have transformed the treatment landscape for
EGFR-mutant NSCLC; however, acquired resistance and the high cost
of therapy remain major clinical challenges, particularly in
resource-limited South Asian settings. This study presents a graph
neural network (GNN)-based computational framework for the
classification of EGFR inhibitors and the prospective virtual
screening of bioactive compounds from traditional Nepali medicinal
plants.

\vspace{0.3cm}
\noindent
A curated dataset of 10,523 unique compounds with experimentally
validated EGFR bioactivity (74\% active) was assembled from ChEMBL,
standardised with RDKit, and represented as molecular graphs using
29-dimensional atom and 12-dimensional bond features. Seven
classification models were benchmarked under a rigorous
Bemis-Murcko scaffold split: five GNN architectures (GCN, GAT, GIN,
AttentiveFP, and GraphSAGE) and two baselines (Random Forest and a
multi-layer perceptron on 2048-bit ECFP4 fingerprints). The
fingerprint baselines led on the scaffold test set (Random Forest
AUROC 0.8827, MLP 0.8547), ahead of the graph neural networks
(0.8087--0.8454, GCN strongest); Bayesian hyperparameter
optimisation with Optuna raised AttentiveFP's validation AUROC to
0.8632 but it did not surpass the baselines, indicating that
well-tuned fingerprint models remain highly competitive on a
scaffold-split dataset of this size.

\vspace{0.3cm}
\noindent
The two best-performing models---the best graph network (GCN) and
the best overall classifier (Random Forest)---were then applied to a
virtual screening library of natural product compounds curated from
seven traditional Nepali medicinal plant species
(\textit{Swertia chirayita}, \textit{Rhododendron arboreum},
\textit{Nardostachys jatamansi}, \textit{Terminalia chebula},
\textit{Berberis aristata}, \textit{Ocimum sanctum}, and
\textit{Tinospora cordifolia}) via the COCONUT~2.0 database. Of the
535 retrieved compounds, eight already present in the ChEMBL training
data---including the apparent top hits apigenin and kaempferol---were
excluded as leakage, leaving 527 genuinely unseen compounds. The
screening hit list proved strongly model-dependent (Random Forest 179
predicted actives versus GCN 24), and on the de-duplicated library
\emph{neither} best model produced a Lipinski-compliant lead at a
predicted probability $\geq 0.80$. The strongest prospective
candidates are the xanthone \textbf{isogentisin} (Random Forest 0.728)
and the flavone \textbf{luteolin}, the top cross-model consensus lead
(both models $\geq 0.50$)---both at moderate confidence.

\vspace{0.3cm}
\noindent
This work shows that model-consensus virtual screening can surface
drug-like natural-product candidates from Nepali ethnomedicinal
libraries using an entirely open-source, CPU-only pipeline, while
underscoring that such screens are strongly model-dependent, must
exclude compounds already seen in training to remain genuinely
prospective, and yield only moderate-confidence leads requiring
experimental validation.
\end{onehalfspace}

\vspace{0.5cm}
\noindent{\fontsize{10}{12}\selectfont\textit{\textbf{Keywords:}
EGFR inhibitors, graph neural networks, virtual screening,
Nepali medicinal plants, drug discovery.}}